\section{\gai Benchmark Design}
\label{section:43-GAI-Design}

Our design of the \gai benchmarking framework draws jointly from the established methods of benchmark design, methods for the empirical study of algorithms~\cite{Hooker1994, Hooker1995} and the theory of constraints~\cite{Goldratt1990, Goldratt2004}.
The prior work in benchmark design guides \textit{what} components to include in the \gai benchmark, namely the \textit{datasets, tasks, metrics} and \textit{system under test}.
The prior work in the empirical analysis of algorithms determines \textit{how} we compose those components to measure the effectiveness and efficiency of a \ac{SUT} for a specific workload.
The prior work in constraint theory determines how we identify bottlenecks and design inputs to trigger bottleneck behavior in a \ac{SUT}.

\subsection{Empirical Benchmark Design}

We elect to design our \gai benchmark empirically, which involves applying the scientific method to drive the selection of datasets, tasks, and metrics that are included in the benchmark. 
The empirical process provides a formal structure that enables transparent, reproducible benchmark development, which we describe in this section.
The empirical benchmark design process is marked by the iterative development and testing of hypotheses that map from input datasets to expected outputs. 
The hypothesis-based mapping allows a benchmark user to infer the system's internal state based on the observed (input, output) pair.

Our approach draws from John Hooker's framework for the empirical analysis of algorithms, summarized in Figure~\ref{Figure:gai-empirical-method}. 
His method identifies four key components: the \textit{problem}, the \textit{approach}, the \textit{phenomena}, and the \textit{apparatus}.
Before explaining the method, we define these four components.
The \textit{problem} here studies what we are trying to do (achieve some outcome in the domain and problem space).
The \textit{approach} here studies the dataset and task features we can use to stimulate the problem (identify datasets, tasks).
The \textit{phenomena} studies the interactions of an approach with a problem (identify suitable metrics and where to find them).
The \textit{apparatus} are the tools and techniques we use to investigate, measure, or otherwise study those interactions between the phenomena and the problem (get datasets, implement tasks, instrument system, monitor execution, etc.).

    
    \begin{figure}[h!]
        \begin{center}
            \includesvg[width=0.5\columnwidth]{Figures/SVGs/empirical_method}
            \caption{Hooker's methodology for an empirical science of algorithms~\cite{Hooker1994,Hooker1995}.}
            \label{Figure:gai-empirical-method}
        \end{center}
    \end{figure}



The process begins by \textbf{Identifying and Characterizing a Relevant Problem}, selecting a use case to study, and identifying the typical tasks, datasets, and challenges. 
\textbf{Identifying and Characterizing the Phenomena of Interest} then investigates the `interesting things' we expect to observe and explains why measuring them helps us understand the underlying problem.
After characterizing the problem and phenomena \textbf{Exploratory Experimentation} gathers the foundational understanding of the bottlenecks, friction points, curiosities, and insights about how to stimulate those bottlenecks.
These insights are used to design meaningful experiments and valid \textbf{Hypotheses}.
The hypotheses focus the analysis and inform data collection by specifying what phenomena we expect to observe given a problem and an approach under study.
With the hypotheses established and the phenomena identified, the next step is to \textbf{Construct the Investigation Apparatus} that will coordinate interactions between the benchmark and the system, generate tasks and datasets, and collect the observable behaviors to produce the metrics of interest.
Finally, the experimenter must \textbf{Analyze the Results, Develop a Theoretic Explanation and Iterate} to move closer to understanding the foundational truths of the problem being studied.

\subsection{Design using the Theory of Constraints}

The design of the \gai framework fuses Hooker's empirical approach with Goldratt's constraint theory~\cite{Goldratt1990, Goldratt2004}.
Constraint theory uses `five focusing steps' to identify and remediate system bottlenecks. 
First, \textbf{identify} the bottleneck, then \textbf{exlploit} it by improving the system's efficiency surrounding the bottleneck.
Next, \textbf{subordinating} the rest of the system to the bottleneck sets the system's throughput equal to that of the bottleneck to prevent a buildup of work-in-progress.
Eventually, we aim to \textbf{elevate} the bottleneck by expanding its capacity until it is no longer the bottleneck before \textbf{iterating} to find the new constraint. 
We leverage this process in reverse order by designing tasks and datasets that promote potential bottlenecks to constrain the system's performance. 
It is then up to the \ac{SUT} to identify how to use the acceleration techniques to exploit or elevate the constraints the adversarially shaped input data causes for the problem-centric tasks.

% Datasets should be drawn from the \textit{Mesh} and \textit{Scale-Free} distributions~\cite{Beamer2015}.

% Roofline Analysis is a useful evaluation tool to understand where the bottleneck in execution is~\cite{Williams2009,Hennessy2019}

% \subsection{Collecting Metrics}

% google ~\cite{Ren2010,Kanev2015,Hennessy2019}


% \osullikomment{Moved from related work.}
% \subsection{Accuracy Bottlenecks}
% \nrscomment{this reads like something youre inventing out of thin air. its not really a bottleneck in the sense of being the limiting reagent. I think a better place to put this (and jsut call it the law of diminishing return) is in the benchmark design considerations in ch4}
% What we call an accuracy bottleneck occurs when the desired degree of accuracy (or some other effectiveness) metric is beyond the elbow point for diminishing returns, where the rate of improvement levels off proportional to the time and compute applied to it. 
% Accuracy bottlenecks are not well explored in \gai benchmarking literature but are a well-known trade-off in machine learning more broadly.
% The trade-off is recognized at the macro-level with the recognition that the amount of computing power invested into deep learning is yielding ever smaller returns~\cite{Thompson2021} and at the micro level avoiding overfitting has relevant strategies like early stopping~\cite{Ying2019} that can be employed.
% Detecting and acting on accuracy bottlenecks is important if a \gai benchmark is going to be capable of trading off between effectiveness, efficiency and cost. 


% \osullikomment{Batch Latency as the metric for all ET?}